% NEW 2026-09-18 (§5.2 readability pass; handoff 2026-09-17_ch5_s52_setup_critique.md
% §V4, §W). REPLACES tab_5-2c_measures, which is deleted. Three changes, each a
% ruling of Marc's from the 2026-09-18 read.
% (1) MEASURES -> METRICS. The chapter's own structure is already built on the
%     word: §5.4 and §5.5 are clustered by instrumented metric family, and
%     Table~\ref{tab:mtd-metrics} is "Metric families in MTD evaluation".
% (2) THE FAMILY COLUMN IS THE TIE-BACK TO THE LITERATURE REVIEW, and it is
%     literal rather than asserted: the cells are Table 3.1's own inner key
%     (success events, attacker time, resource spent, system state,
%     network-state change), so a reader can put every metric this chapter
%     reports back on the field's map and see that what judges a defence here is
%     the field's measure and not this thesis's. Marc's ask: "consider what
%     metrics we can tie back into the metric table that we produced in the
%     literature review."
% (3) THE FOUR §5.3 INSTRUMENTS ARE GONE from this table --- campaign coverage,
%     opening variety, profile divergence, response to disruption. Marc's ruling
%     on the self-scoring question: calling them validation is a claim of proof a
%     self-chosen instrument set cannot carry ("you can be selective about it ---
%     that's the whole issue"). Evaluation conventions §f has the honest word:
%     attacker-property claims are taken from the no-defence arm and reported as
%     OBSERVATIONS, and the verdict they add up to belongs in the discussion. They
%     are therefore declared in §5.3, where they are first used (conventions §a),
%     and this table holds only what evaluates a defence. Twelve rows became eight.
% THE ONE DECLARED DIVERGENCE: deployment tempo is MTD execution frequency, which
% Table 3.1 files under EFFECTIVENESS (network-state change, defender-side); this
% chapter reports it in §5.5 as a cost. The rows are grouped by the section that
% reports the metric, and the caption says so, so the divergence is declared
% rather than hidden. Flagged for Marc: the alternative is to move the row to the
% effectiveness group and report it there.
% EVERY ROW IS A QUANTITY THE CORPUS ACTUALLY PRODUCES: transcribed from
% data/results/ch5_defended/analyse.py (summarise_movement, summarise_baseline).
% A metric that is not computed does not appear --- internal MTTC stays out until
% its brief is ruled, and the compromise-checkpoint metric until the run-length
% ruling takes it. Blocked fraction is structurally zero on the baseline arm
% (analyse.py:166), which is why its comparability cell says movement arm.
% CAPTION SCOPED 2026-09-18: "every metric this chapter reports" was false ---
% §5.3's four instruments are metrics this chapter reports and are deliberately
% not here (see the ruling above). "Judges a defence by" is the true scope and
% is the phrase §5.2's prose now uses, so caption and prose agree.
% OPEN FOR MARC 2026-09-18 --- DO WE CARRY THE LINEAGE'S METRIC NAMES? He asked
% of host breadth, "is that HCR, network compromise ratio? Are we not carrying
% over the names of our lineage's metrics?" Checked, and the answer differs per
% row.
%  HOST BREADTH IS NOT HCR/NCR, on two counts. (i) It is not the same quantity:
%   HCR is Ho's Eq. 10, C_t / T_host at a compromise CHECKPOINT, bounded [0,1]
%   (evaluation.py:126, compromised_num / host_num), and Zhang's NCR is the
%   checkpoint metric that ends the run and fixes where MTTC is read; ours is the
%   distinct hosts reached across a whole run, and the number actually REPORTED
%   is its suppression against the no-defence arm (analyse.py:227). On a fixed
%   50-host network a count and a ratio differ by a constant, but a difference
%   from a control is neither. (ii) NEITHER PAPER REPORTS IT AS A RESULT: Ho
%   declares eleven selector FEATURES and four EVALUATION metrics (ASR, RoA, APE,
%   Risk) and HCR is a feature (ho2024.md L86); Tay reports the same five minus
%   RoA's sibling. Carrying the name would import a label the lineage does not
%   use for an outcome. The tie-back that IS true is already in this table: the
%   Family cell "system state" is Table 3.1's row holding HCR and NCR.
%  WHERE WE DO DROP A LINEAGE NAME, AND IT IS WORTH HIS RULING: Brown's two
%   metrics are "attack actions blocked" (Fig. 4) and "the average attempts
%   required to compromise" (Fig. 5) --- both carried under Brown's own names in
%   Table 3.1 (:2404, :2425) --- and they are this table's BLOCKED FRACTION and
%   EFFORT PER HOST, renormalised (a share rather than a total; per distinct host
%   rather than per compromise). Brown is the direct substrate lineage and §5.4.3
%   re-runs his claim, so different names for his own two metrics cost something
%   and buy nothing.
%  THE PRECEDENT FOR THE MIDDLE PATH IS THIS REPO'S OWN: "internal MTTC" keeps
%   the lineage name, qualifies it in one word, and writes the divergence down
%   (metrics_semantics.md §a, §c). "Attack actions blocked, as a share" would be
%   the same move. NOT DONE HERE --- renaming reaches the generated §5.4/§5.5
%   captions, tab:eff-lineage and §5.4's prose, which is Marc's.
% Geometry tested against the real class and geometry: no overfull box; the
% tabular measures 454.1 pt against a 455.2 pt text width, body height 283 pt.
% SHOULD BE BORN GENERATED, with tab_5-2a, by the reader that produces the
% §5.3--§5.5 floats (C39). Ratify on read.
\begin{table}[htbp]
  \centering
  \caption[The metrics]{Every metric this chapter judges a defence by, the family of Table~\ref{tab:mtd-metrics} it instantiates, and what it may be compared across. The rows are grouped by the section that reports the metric rather than by that table's own filing, which places deployment tempo under effectiveness. Numbers published on other simulators are comparable to none of these.}
  \label{tab:metrics}
  \tablestyle\setlength{\tabcolsep}{4pt}
  \begin{tabular}{@{}cP{2.0cm}P{2.8cm}P{5.0cm}P{2.7cm}@{}}
    \toprule
    & Family & Metric & What it is & Comparable across \\
    \midrule
    & system state & Host breadth & the distinct hosts compromised in a run, and its suppression against no defence & both arms, as a count \\
    & success events & Target reach & whether the run reached the database hosts & both arms, as a fraction \\
    & attacker time & Delay to first compromise & the time to the first host compromised, censored at the run length & within an arm \\
    \rowgroup{4}{Effectiveness} & success events & Blocked fraction & the share of the attacker's actions a deployment interrupts & the movement arm \\
    \midrule
    & resource spent & Attacker cost & the actions attempted, and the run's time split into dwell, imposed delay and the remainder & counts both arms; the time split within the movement arm \\
    & resource spent & Effort per host & the actions attempted for each distinct host compromised & both arms, as a ratio \\
    & resource spent & Reconfiguration occupancy & the share of the run during which a mechanism is deploying & both arms \\
    \multirow{-4}{*}[0.18cm]{\smash{\rotatebox{90}{\emph{Efficiency}}}} & network-state change & Deployment tempo & the deployments executed per thousand seconds & both arms \\
    \bottomrule
  \end{tabular}
\end{table}
